\documentclass[11pt]{amsart}
\usepackage[a4paper,margin=1in]{geometry}
\raggedbottom
\usepackage[bottom]{footmisc}
\usepackage{amsmath,amssymb,amsthm}
\usepackage{booktabs}
\usepackage[hidelinks]{hyperref}
\hypersetup{pdftitle={A local non-containment in Step (v) of IUT IV, Theorem 1.10},
  pdfauthor={Toshihisa Urahigashi}}

\newtheorem{theorem}{Theorem}
\newtheorem{proposition}[theorem]{Proposition}
\newtheorem{lemma}[theorem]{Lemma}
\theoremstyle{remark}
\newtheorem{remark}[theorem]{Remark}

\newcommand{\Q}{\mathbb Q}
\newcommand{\Z}{\mathbb Z}
\newcommand{\F}{\mathbb F}
\newcommand{\Qp}{\mathbb Q_p}
\newcommand{\Zp}{\mathbb Z_p}
\newcommand{\cO}{\mathcal O}
\newcommand{\cA}{\mathcal A}
\newcommand{\cI}{\mathcal I}
\newcommand{\cU}{\mathcal U}
\newcommand{\dd}{\mathfrak d}
\newcommand{\qq}{\underline{\underline{q}}{}}
\newcommand{\Fmod}{F_{\mathrm{mod}}}
\newcommand{\VV}{\underline{\mathbb V}}
\newcommand{\Vbadmod}{\mathbb V^{\mathrm{bad}}_{\mathrm{mod}}}
\newcommand{\Vbad}{\underline{\mathbb V}^{\mathrm{bad}}}
\newcommand{\mulog}{\mu^{\log}}
\newcommand{\RIt}{(R_I)^{\sim}}
\newcommand{\pr}{\operatorname{pr}}
\newcommand{\Nm}{\operatorname{N}}
\newcommand{\Tr}{\operatorname{Tr}}
\newcommand{\Sj}{\mathbb S^{\pm}_{j+1}}
\newcommand{\Dp}{\mathcal D^{\vdash}_{p}}
\newcommand{\bv}{\mathbf v}
\newcommand{\bw}{\mathbf w}
\newcommand{\bu}{\mathbf u}

\begin{document}

\title{A local non-containment in Step~(v) of IUT~IV, Theorem~1.10}
\author{Toshihisa Urahigashi}
\email{123026.urahigashi.toshihisa@gmail.com}
\date{Draft, 5 October 2026}

\begin{abstract}
Step~(v) of the proof of Theorem~1.10 in the fourth paper of Mochizuki's series on inter-universal
Teichm\"uller theory estimates, prime by prime, the log-volume of the holomorphic hull of the union of
possible images of a $\Theta$-pilot object defined in Corollary~3.12 of the third paper. For each tuple of
places it bounds the log-volume of the component of this hull by that of an explicit module, which depends on the
order of the $q$-parameter at the place carrying the label $j$. We give an explicit initial $\Theta$-datum, with
two places of multiplicative reduction above the prime $30965771$, and an element of the projection of the union
to one summand, obtained from a theta value by a permutation of labels allowed by the indeterminacy (Ind1), that
does not lie in this module; the corresponding component of the hull has larger log-volume than the bound obtained
in Step~(v).  No assertion is made here about Corollary 3.12 itself, the final global inequality of Theorem 1.10, or the
abc conjecture.
\end{abstract}

\maketitle

\section{The statement of Step (v)}\label{s:source}

We use the notation of \cite{IUT1,IUT3,IUT4}; page numbers refer to the versions listed at the end, and we
omit some decorations such as ${}^{n,\circ}$. For a finite extension $k$ of $\Qp$ we write $\cO_k$ for its ring
of integers, $e_k$ for its ramification index, $\operatorname{ord}$ for the valuation on an algebraic closure of
$\Qp$ with $\operatorname{ord}(p)=1$, and $\dd_k$ for the order of a generator of the different of $\cO_k$ over
$\Zp$. For finitely many such fields $(k_i)_{i\in I}$ we write, as in \cite[Props.~1.1, 1.2, pp.~9--10]{IUT4},
\[
k_I=\textstyle\bigotimes_{i\in I}k_i,\qquad R_I=\bigotimes_{i\in I}\cO_{k_i},\qquad
\log_p(R_I^\times)=\bigotimes_{i\in I}\log_p(\cO_{k_i}^\times),\qquad \dd_I=\sum_{i\in I}\dd_{k_i},
\]
the tensor products being over $\Qp$, respectively $\Zp$, and $\RIt$ for the normalization of $R_I$. The ring
$k_I$ is a finite product of fields $L$, and $\RIt$ is the product of their rings of integers. We use the
log-volume $\mulog$ of \cite[Prop.~1.4(i), p.~13]{IUT4}, normalized by $\mulog(\RIt)=0$ and
$\mulog(p\cdot\RIt)=-\log p$, and three elementary facts.
\begin{enumerate}
\item[(1.1)] In every factor $L$ of $k_I$, the component of a pure tensor $y=\otimes_i z_i$, $z_i\ne0$, has
  $\operatorname{ord}$ equal to $c=\sum_i\operatorname{ord}(z_i)$. Hence
  $\RIt\cdot y=\{z\in k_I:\operatorname{ord}(z_L)\ge c\text{ for every }L\}$ and $\mulog(\RIt\cdot y)=-c\log p$.
\item[(1.2)] If $p$ is odd and $e_k\le p-2$, then $\log_p(\cO_k^\times)=\{z\in k:\operatorname{ord}(z)\ge 1/e_k\}$
  \cite[Prop.~1.2(i), p.~10]{IUT4}; so for $z\in k^\times$ and $n\in\Z$ we have
  $z\in p^n\log_p(\cO_k^\times)$ if and only if $n<\operatorname{ord}(z)$.
\item[(1.3)] Write $\cA(k)$ for the group of $\Qp$-linear automorphisms of $k$ that map
  $\log_p(\cO_k^\times)$ onto itself, equivalently the log-shell $p^{-1}\log_p(\cO_k^\times)$ ($p$ odd)
  \cite[Rem.~1.2.2(i), p.~36]{IUT3}. Under the hypothesis of (1.2), every $\psi\in\cA(k)$ maps each
  $p^n\log_p(\cO_k^\times)$ onto itself, hence $\lceil\operatorname{ord}\psi(z)\rceil=\lceil\operatorname{ord} z\rceil$;
  in particular $\operatorname{ord}\psi(z)\le\lceil\operatorname{ord} z\rceil$.
\end{enumerate}

Fix an initial $\Theta$-datum \cite[Def.~3.1, pp.~61--63]{IUT1} that satisfies the hypotheses of
\cite[Thm.~1.10, pp.~22--23]{IUT4}, a prime number $p$ that ramifies in $K$, and $j\in\{1,\dots,\ell^*\}$, where
$\ell^*=(\ell-1)/2$; write $\Sj=\{0,1,\dots,j\}$ \cite[p.~27]{IUT4}. For a place $v$ of $\Fmod$ above $p$ let $\underline v\in\VV$ be the corresponding element
\cite[Def.~3.1(e)]{IUT1}; we write $k_v=K_{\underline v}$, $\qq_v=\qq_{\underline v}$ for the $2\ell$-th root of
the $q$-parameter \cite[Ex.~3.2(iv), p.~71]{IUT1}, and $\rho(v)=\operatorname{ord}_p(\qq_v)$. The packet
$\cI^{\Q}(\Sj;\Dp)$ is the tensor product, over the labels in $\Sj$, of the direct sum over the places $v\mid p$ of
the $\Q$-spans of the log-shells \cite[p.~27]{IUT4}, \cite[Prop.~3.2, pp.~97--98]{IUT3}. Identifying the $\Q$-span
of the log-shell at $v$ with $k_v$ \cite[Rem.~3.1.1(i), p.~93; Prop.~3.2(i)]{IUT3}, we get
\[
\cI^{\Q}(\Sj;\Dp)=\bigoplus_{\bv}k_{\bv},\qquad k_{\bv}=k_{v_0}\otimes k_{v_1}\otimes\cdots\otimes k_{v_j},
\]
the sum running over the tuples $\bv=(v_0,\dots,v_j)$ of places of $\Fmod$ above $p$; the factor of $k_{\bv}$ with
label $i$ is $k_{v_i}$. Write $\pr_{\bv}$ for the projection to $k_{\bv}$, $U_j$ for the union of the possible
images of a $\Theta$-pilot object, subject to the indeterminacies (Ind1), (Ind2), (Ind3), that is formed in the
proof of \cite[Cor.~3.12]{IUT3} (denoted there $\cU_{j,v_{\Q}}$, with $v_{\Q}=p$; p.~174), and $\overline U_j$ for its
holomorphic hull \cite[Rem.~3.9.5(i), p.~127]{IUT3}.

For a tuple $\bv$, Step~(v) \cite[pp.~27--28]{IUT4} applies \cite[Prop.~1.4(iii), p.~13]{IUT4} with
$I=\Sj$, $k_i=k_{v_i}$, $I^*=\{i:e_{k_{v_i}}>p-2\}$, $i^\dagger=j$, and $\lambda=\operatorname{ord}(\qq_{v_j}^{\,j^2})
=j^2\rho(v_j)$ if $\underline v_j\in\Vbad$ ($\lambda=0$ otherwise). From the
inclusion $\phi(p^\lambda\cdot\RIt)\subseteq C_{\bv}$ of that proposition, where
\[
C_{\bv}=p^{\lfloor\lambda-\dd_I-a_I\rfloor}\cdot\log_p(R_I^\times)\ \subseteq\
M_{\bv}=p^{\lfloor\lambda-\dd_I-a_I\rfloor-b_I}\cdot\RIt
\]
and $a_I,b_I$ are as in \cite[Prop.~1.2]{IUT4}, it concludes that a suitable multiple of $\log_p(R_I^\times)$
contains the union of possible images,\footnote{In the words of \cite[p.~27]{IUT4}, the inclusion ``implies that
the result of multiplying $p^{\lfloor\lambda\rfloor-|I|}\cdot2^{-|I|}\cdot\log_p(R_I^\times)$ by a suitable nonpositive
[\dots] integer power of $p_{v_{\Q}}$ contains the `union of possible images of a $\Theta$-pilot object' discussed in
Step~(iv)''. Which power is meant plays no role below: Theorem~\ref{t:main} concerns $M_{\bv}$ and the
bound~\eqref{eq:bound}.} and that an upper bound for the component at $\bv$ of the log-volume of the holomorphic
hull may be obtained from an upper bound for the log-volume of $M_{\bv}$ \cite[p.~28]{IUT4}. If $I^*=\emptyset$, the
latter is
\begin{equation}\label{eq:bound}
\mulog(M_{\bv})\le(-\lambda+\dd_I+1)\cdot\log p
\end{equation}
\cite[Prop.~1.4(iii)]{IUT4}. This is the bound that we test.

\begin{lemma}\label{l:object}
The union of possible images whose holomorphic hull Step~(v) estimates is $U_j$, and its component at the
tuple $\bv$ is $\pr_{\bv}(U_j)$.
\end{lemma}

\begin{proof}
Step~(v) states that the container ``contains the `union of possible images of a $\Theta$-pilot object'
discussed in Step~(iv)'' \cite[p.~27]{IUT4}. Step~(iv) introduces this union ``according to the various
definitions given in the statement of [IUTchIII], Corollary~3.12'' \cite[p.~26]{IUT4}, and estimates, for fixed
$v_{\Q}$ and $j$, its component in $\cI^{\Q}(\Sj;\mathcal D^{\vdash}_{v_{\Q}})$ and then the component
corresponding to each tuple $\{v_i\}$ \cite[p.~27]{IUT4}. In the proof of \cite[Cor.~3.12]{IUT3} this union is the
subset $\cU_{j,v_{\Q}}\subseteq\cI^{\Q}(\Sj;\mathcal D^{\vdash}_{v_{\Q}})$, formed before its holomorphic hull is
taken \cite[p.~174]{IUT3}. For $v_{\Q}=p$, its component at the tuple $\bv$ is $\pr_{\bv}(U_j)$.
\end{proof}

\section{The datum \texorpdfstring{$D_2$}{D2}}\label{s:datum}

Let $F_0=\Q(\sqrt5)$, $\ell=7$, and
\[
\pi=14164+5825\sqrt5,\qquad p=\Nm_{F_0/\Q}(\pi)=30965771,\qquad a=\pi\bar\pi^{\,6},
\]
where the bar denotes the conjugation of $F_0$. Let $E_0\colon y^2=x(x-1)(x-a)$ over $F_0$,
$F=F_0(\sqrt{-1},E_0[15])$, $K=F(E_0[7])$, $X_F$ the curve $E_0\times_{F_0}F$ minus its origin, and $\Vbadmod$ the
set of places $v$ of $F_0$ with $v\nmid 14$ and $\operatorname{ord}_v(j(E_0))<0$.

\begin{proposition}\label{p:datum}
These data extend to an initial $\Theta$-datum
$D_2=(\overline F/F,X_F,\ell,\underline C_K,\VV,\Vbadmod,\underline\epsilon)$ in the sense of
\cite[Def.~3.1]{IUT1} that satisfies the hypotheses of \cite[Thm.~1.10]{IUT4}, with $\Fmod=F_0$. The number $p$ is
prime and splits in $F_0$; the two places $b_1=v_\pi$ and $b_2=v_{\bar\pi}$ above $p$ lie in $\Vbadmod$, have
$[(\Fmod)_{b_i}:\Qp]=1$, and have the local data of Table~\ref{tab:D2}. Moreover
$p>e^*_{\mathrm{mod}}\cdot\ell=7741440$.
\end{proposition}

\begin{table}[ht]
\centering
\begin{tabular}{@{}lcccccc@{}}
\toprule
place $v$ & $\operatorname{ord}_v(a)$ & $\operatorname{ord}_p(q_v)$ & $e_{k_v}$ & different exponent & $\dd_{k_v}$
  & $\rho(v)$\\
\midrule
$b_1$ & $1$ & $2$ & $105$ & $104$ & $104/105$ & $1/7$\\
$b_2$ & $6$ & $12$ & $35$ & $34$ & $34/35$ & $6/7$\\
\bottomrule
\end{tabular}
\caption{The places of $D_2$ above $p=30965771$; $k_v=K_{\underline v}$ and $\rho(v)=\operatorname{ord}_p(\qq_v)$.}\label{tab:D2}
\end{table}

\begin{proof}[Proof of Proposition~\ref{p:datum}, except Table~\ref{tab:D2}]
The integer $p$ is prime, $p\equiv1\pmod5$ and $p\equiv3\pmod4$. So $p=\pi\bar\pi$ splits in $F_0$, both places
above it have $[(F_0)_{b_i}:\Qp]=1$, and $\Qp(\sqrt{-1})$ is unramified over $\Qp$.

\emph{The integer $\Nm(1-a)$.} We have $\Nm(1-a)=p^7-p\Tr(\pi^5)+1=2^2\cdot11^2\cdot R$, where
\[
R=56406340220143710508313748934073550973772664592051
\]
is prime. Hence no prime divides $\Nm(1-a)$ to the seventh power, $\operatorname{ord}_v(1-a)\le6$ at every place $v$ of $F_0$,
and $\operatorname{ord}_2(1-a)=1$ at the place above $2$, which is inert. Neither $\Nm(a)=p^7$ nor $\Nm(1-a)$ is
divisible by $7$.

\emph{The curve.} We have $j=j(E_0)=2^8(a^2-a+1)^3/(a^2(a-1)^2)$ and $a\in\cO_{F_0}$. At an odd place $v$ with
$\operatorname{ord}_v(a)>0$ or $\operatorname{ord}_v(1-a)>0$ the element $a^2-a+1$ is a unit, so
$\operatorname{ord}_v(j)=-2\operatorname{ord}_v(a)$, respectively $-2\operatorname{ord}_v(1-a)$; at the other odd
places the discriminant $16a^2(a-1)^2$ is a unit. At $b_1$ and $b_2$, $1-a$ is a unit, so
$\operatorname{ord}_{b_1}(j)=-2$ and $\operatorname{ord}_{b_2}(j)=-12$. A rational $j$ would have equal orders at the
two places above $p$; hence $j\notin\Q$ and $\Fmod=\Q(j)=F_0$ \cite[Prop.~1.8(ii), p.~19]{IUT4}. At the place above
$2$, $a$ is a unit and $\operatorname{ord}_2(j)=8-2=6\ge0$.

\emph{Conditions (a), (b) of \cite[Def.~3.1]{IUT1}.} We have $\sqrt{-1}\in F$. As $E_0[2]$ is rational over $F_0$,
$F=F_0(\sqrt{-1},E_0[30])$; it is Galois over $F_0$, its Galois group embeds in $\Z/2\times\mathrm{GL}_2(\Z/15\Z)$,
whose order is $2^{10}\cdot3^2\cdot5$, so $[F:F_0]$ is prime to $7$; and $E_F[6]$ is rational over $F$. By \cite[Prop.~1.8(v), p.~19]{IUT4},
applied with the primes $3$ and $5$, $E_F$ has semistable reduction at every finite place of $F$; so its reduction
is multiplicative where $\operatorname{ord}(j)<0$ and good where $\operatorname{ord}(j)\ge0$
\cite[Props.~VII.5.1, VII.5.5]{Sil}. By the previous paragraph, the places of multiplicative reduction are exactly
the places above $\Vbadmod$; this set is nonempty ($b_1,b_2\in\Vbadmod$), its residue characteristics are odd, and
$E_F$ has good reduction at every other place not dividing $2\cdot7$, which is the reduction hypothesis of
\cite[Thm.~1.10]{IUT4}. The field $F$ has the form required there: $F_{\mathrm{tpd}}=\Fmod(E[2])=F_0$, $E_0$ is in
Legendre form, and $F=F_{\mathrm{tpd}}(\sqrt{-1},E_0[3\cdot5])$.

\emph{Condition (c).} The prime $\ell=7$ is not a residue characteristic of $\Vbadmod$. At a place $w$ of $F$ above
$v\in\Vbadmod$ the $q$-parameter has order $-\operatorname{ord}_w(j)=2\,e(w|v)\,r$, where
$r=\operatorname{ord}_v(a)$ or $\operatorname{ord}_v(1-a)$ lies in $\{1,\dots,6\}$ ($a$ vanishes only at $b_1$ and
$b_2$, to the orders $1$ and $6$) and $e(w|v)$ divides $[F:F_0]$; so these orders are prime to $7$. At such a
place $w$ the Legendre equation is minimal, as $c_4=16(a^2-a+1)$ is a unit \cite[\S VII.1]{Sil}, and its reduction
has a node, at $(0,0)$ with tangents $y=\pm\sqrt{-1}\,x$ if $a\equiv0$ and at $(1,0)$ with tangents $y=\pm(x-1)$ if
$a\equiv1$, rational over the residue field; so $E_F$ has split multiplicative reduction and $E_F\cong E_q$ over
$F_w$ \cite[App.~C, Thm.~14.1]{Sil}. For the image of $G_F$ in $\mathrm{GL}(E_0[7])\cong\mathrm{GL}_2(\F_7)$: above
$b_1$, $a\equiv0$ and $\operatorname{ord}_w(q)=2e(w|b_1)$ is prime to $7$. As
$q^{1/7}$ generates a ramified extension of the maximal unramified extension of $F_w$, which contains $\mu_7$, an
element of inertia acts on $E_0[7]\cong\langle\zeta_7,q^{1/7}\rangle$ by a nontrivial transvection $t$. At the
place of $F_0$ of degree one above $19$ at which $\sqrt5\equiv9$, $E_0$ has good reduction, $a\equiv15$ and
$a_{19}=19+1-\#\tilde E_0(\F_{19})=-4$; the Frobenius $g$ acts on $E_0[7]\cong\tilde E_0[7]$
\cite[Prop.~VII.3.1]{Sil} and satisfies $g^2+4g+19=0$ \cite[Thm.~V.2.3.1]{Sil}; as $X^2+4X+5$ is irreducible over
$\F_7$ (its discriminant $3$ is not a square), $g$ fixes no line. The image $H$ of $G_F$ is normal in the image of
$G_{F_0}$, so it contains $t$ and $gtg^{-1}$, transvections with distinct fixed lines, which generate
$\mathrm{SL}_2(\F_7)$.

\emph{Conditions (d)--(f).} We do not invoke \cite[Cor.~2.2]{IUT4} for $D_2$: its condition (P1)
\cite[p.~45]{IUT4} contains the inequality $h^{1/2}\le\ell$ for the height $h$ defined in its proof (p.~44), and here
$h>\ell^2$. We construct the data directly; as noted on p.~46, the essential input for $\VV$ and
$\underline\epsilon$ is the property, shown above, that the image of $G_F$ contains $\mathrm{SL}_2(\F_\ell)$. Since
$j\notin\Q$, $j$ is none of $2^{14}\cdot31^{3}\cdot5^{-3}$, $2^{2}\cdot73^{3}\cdot3^{-4}$, $1728$, $0$, so $X_F$ admits an
$F$-core \cite[Prop.~2.1]{MFHMP}. Hence $X_{\overline F}$ is not arithmetic \cite[Prop.~2.3(i)]{CanLift}, nor is its
finite \'etale quotient $C_{\overline F}=[X_{\overline F}/\{\pm1\}]$, a punctured hemi-elliptic orbicurve, which is
therefore an $\overline F$-core \cite[Prop.~2.7]{CanLift}; so $C_F=[X_F/\{\pm1\}]$ and $C_K=C_F\times_FK$ are an
$F$-core and a $K$-core \cite[Prop.~2.3(i)]{CanLift}. The group $E_0[7]$ is rational over $K$. Fix a line
$W\subset E_0[7]$ and $0\ne\xi\in E_0[7]/W$; let $\alpha\colon E_K\to E'=E_K/W$ be the quotient and
$\beta\colon E'\to E_K$ the isogeny with $\beta\alpha=[7]$. The finite \'etale cover
$\beta\colon\underline X_K=E'\setminus\ker\beta\to X_K$ has group $\ker\beta=\alpha(E_0[7])\cong E_0[7]/W$,
corresponds to the quotient $E_0[7]\to E_0[7]/W$ of the maximal abelian quotient modulo~$7$ of the geometric
fundamental group of $X_K$, and has its cusps, the points of $\ker\beta$, rational over $K$; so
$\underline C_K=[\underline X_K/\{\pm1\}]$, formed with the cusp $0$ as origin, is of type $(1,7\text{-tors})^{\pm}$
\cite[Def.~2.1]{EtTh}, and, being finite \'etale over the $K$-core $C_K$, it has $K$-core $C_K$, as required in (d).
Let $\underline\epsilon$ be the cusp of $\underline C_K$ given by $\pm\xi$, as in (f). At a place $w$ of $K$ above an
element of $\Vbadmod$, the Tate uniformization of $E_0$ over $K_w$, which exists by the paragraph on (c), gives a line
$\mu_7\subset E_0[7]$ and a generator $\pm1$, canonical up to sign, of $E_0[7]/\mu_7\cong\Z/7\Z$; if $(W,\pm\xi)$ is
this pair, then $\underline C_K$ is of type $(1,\Z/7\Z)^{\pm}$ at $w$ \cite[Def.~2.5(i)]{EtTh} and
$\underline\epsilon_w$ is the cusp of $\pm1$, as required in (e), (f). The group $\mathrm{SL}_2(\F_7)$ acts
transitively on the pairs $(L,\pm\eta)$ of a line $L\subset E_0[7]$ and $0\ne\eta\in E_0[7]/L$ (it sends any basis of
determinant $1$ to any other), and it lies in the image of $\mathrm{Gal}(K/F)$. Since $E_0$ and its Tate parameter are
defined over the completion of $F$ below $w$, an element $\gamma\in\mathrm{Gal}(K/F)$ carries the pair at $w$ to the
pair at $\gamma(w)$, and we may choose $\gamma$ so that the latter is $(W,\pm\xi)$. Choosing such a place over each element of
$\Vbadmod$, and any place over the others, gives $\VV$.

\emph{The threshold.} $\Fmod=\Q(\sqrt5)$ has $e_{\mathrm{mod}}=2$, so
$e^*_{\mathrm{mod}}\cdot\ell=2^{12}\cdot3^3\cdot5\cdot2\cdot7=7741440<p$ \cite[p.~22]{IUT4}.
\end{proof}

\section{The places above \texorpdfstring{$p$}{p}}\label{s:local}

\begin{proof}[Proof of Table~\ref{tab:D2}]
Let $b\in\{b_1,b_2\}$ and $m=-\operatorname{ord}_b(j)$, so $m=2$ at $b_1$ and $m=12$ at $b_2$; let
$k_0=(F_0)_b(\sqrt{-1})=\Qp(\sqrt{-1})$. As in the previous proof, $E_0$ has split multiplicative reduction over
$k_0$, so $E_0\cong E_q$ over $k_0$ with $\operatorname{ord}_p(q)=m$, and for $n$ prime to $p$ the field
$k_0(E_0[n])$ is $k_0(\mu_n,q^{1/n})$ \cite[App.~C, Thm.~14.1, and the expansion $j=q^{-1}+744+\cdots$ on
p.~444]{Sil}. Write $q=u\,p^m$ with $u\in\cO_{k_0}^\times$. The field $k_0(\mu_n,u^{1/n})$ is unramified over
$\Qp$, and adjoining to it an $n$-th root of $p^m$ gives a totally ramified extension of degree $n/\gcd(n,m)$.
Hence $k_0(E_0[n])$ is tamely ramified over $\Qp$, with ramification index $n/\gcd(n,m)$. For $\underline v\in\VV$
above $b$ we have $k_v=K_{\underline v}=k_0(E_0[105])$, so $e_{k_v}=105/\gcd(105,m)$, which is $105$ at $b_1$ and
$35$ at $b_2$. Since $k_v$ is tame over $\Qp$, its different has exponent $e_{k_v}-1$, that is
$\dd_{k_v}=(e_{k_v}-1)/e_{k_v}$ \cite[Prop.~1.3(i), pp.~11--12]{IUT4}. Finally
$\qq_v=q^{1/2\ell}$ \cite[Ex.~3.2(iv), p.~71]{IUT1}, so $\rho(v)=m/14$, which is $1/7$ at $b_1$ and $6/7$ at
$b_2$. Both $e_{k_v}$ are at
most $p-2$.
\end{proof}

The integer calculations of \S\S\ref{s:datum}--\ref{s:local} were independently checked by two implementations.

\section{Two elements of the union}\label{s:union}

\begin{lemma}\label{l:member}
Let $\bv$ be a tuple with $\underline v_j\in\Vbad$, and let $\theta\in k_{v_j}$ be a generator of the component
labelled $j$ at $\underline v_j$ of the splitting monoid, that is, the theta value $\qq_{v_j}^{\,j^2}$ up to a root
of unity \cite[p.~11]{IUT3}; so $\operatorname{ord}(\theta)=j^2\rho(v_j)$. Let $x\in k_{\bv}$ be the pure tensor
with $\theta$ at the factor of label $j$ and $1$ at the other factors. Then $U_j$ contains an element $X$ with
$\pr_{\bv}(X)=\chi(x)$ for some $\chi=\otimes_i\chi_i$, $\chi_i\in\cA(k_{v_i})$.
\end{lemma}

\begin{proof}
The $\Theta$-pilot object is determined by generators, at the places of $\Vbad$, of the monoids of
\cite[Def.~3.8(i), p.~112]{IUT3}, and it is constructed from the local fractional ideals generated by elements of
these monoids \cite[Prop.~3.7(v), p.~112]{IUT3}. At $\underline v\in\Vbad$ the component labelled $j$ of these
monoids is a subset of the part of the packet whose factor of label $j$ is the summand at $\underline v$, and it acts
on this part multiplicatively \cite[Prop.~3.4(ii), pp.~102--103]{IUT3}; it lies there through the natural
homomorphism of \cite[Prop.~3.4(i)]{IUT3}, that is, by tensoring with $1$ in the factors of the other labels
\cite[Prop.~3.1(ii), p.~93]{IUT3}. A fractional ideal generated by an element contains that element. The possible
images of \cite[Cor.~3.12]{IUT3} are the images of the $\Theta$-pilot object relative to the Kummer isomorphisms of
\cite[Thm.~3.11(ii), p.~155]{IUT3}, which include isomorphisms of these splitting monoids and of the parts of the
packets just described, in the multiradial representation of \cite[Thm.~3.11(i)]{IUT3}, subject to (Ind1), (Ind2),
(Ind3), and $U_j$ is their union \cite[pp.~173--174]{IUT3}. Let $g$ be the element with $\theta$ in the factor of
label $j$ and $1$ in the other factors, in the packet of the Hodge theater in which the $\Theta$-pilot object is
formed. By the above it lies in the corresponding region of the $\Theta$-pilot object, so $U_j$ contains
$X=\kappa(g)$, where $\kappa$ is the Kummer isomorphism of packets onto $\cI^{\Q}(\Sj;\Dp)$ of
\cite[Thm.~3.11(ii)(a)]{IUT3}, built place by place from isomorphisms of log-shells \cite[Prop.~3.2]{IUT3}. In that
Hodge theater the $\Q$-spans of the log-shells are the fields $k_v$ themselves, and the component of $g$ at $\bv$ is
$x$. For $\cI^{\Q}(\Sj;\Dp)$ the identifications of the $\Q$-spans of the log-shells with the fields $k_v$, used in
\S\ref{s:source}, form an $\mathrm{Ism}$-orbit compatible with the log-shells \cite[Prop.~1.2(vi), p.~32]{IUT3}. Read
through them, $\kappa$ maps each summand $k_{\bw}$ to itself by a tensor product of automorphisms that preserve the
log-shells, so
\[
\pr_{\bv}(X)=\pr_{\bv}(\kappa(g))=\chi\bigl(\pr_{\bv}(g)\bigr)=\chi(x),\qquad\chi=\otimes_i\chi_i,\quad
\chi_i\in\cA(k_{v_i}).
\]
\end{proof}

For a permutation $\sigma$ of $\Sj$ and a tuple $\bv$, let $\sigma_*\colon k_{\bv}\to k_{\bv\circ\sigma^{-1}}$ send the
factor of label $i$ to the factor of label $\sigma(i)$: $\sigma_*(\otimes_i z_i)=\otimes_i z_{\sigma^{-1}(i)}$.

\begin{lemma}\label{l:transport}
Let $\sigma$ be a permutation of $\Sj$ with $\sigma(0)=0$, let $X\in U_j$, let $\bv$ be a tuple, and put
$\bw=\bv\circ\sigma^{-1}$. The images of $X$ under the indeterminacy (Ind1) attached to $\sigma$, followed by (Ind2),
exist and lie in $U_j$, and the component at $\bw$ of each of them has the form $\psi(\sigma_*(\pr_{\bv}X))$ with
$\psi=\otimes_i\psi_i$, $\psi_i\in\cA(k_{w_i})$. In particular $\psi(\sigma_*(\pr_{\bv}X))\in\pr_{\bw}(U_j)$.
\end{lemma}

\begin{proof}
A morphism of processions consists of an order-preserving injection of the stages and, at each stage, a
capsule-full poly-morphism, with no condition relating the stages \cite[Def.~4.10, p.~120]{IUT1}; a capsule-full
poly-isomorphism over a bijection of index sets is the set of all collections of isomorphisms between the
corresponding members \cite[\S0, p.~34]{IUT1}. So $\sigma$ at the stage with index set $\Sj$, together with the
identity at the other stages, is an automorphism of the procession of \cite[Thm.~3.11(i)]{IUT3}, whose index sets
the source regards as ``subject to arbitrary permutation automorphisms'' \cite[p.~10]{IUT3}; cf.\
\cite[Prop.~6.9(i), p.~169]{IUT1}. (Ind1) consists of the indeterminacies induced by such automorphisms, and
(Ind2) of the action of independent copies of $\mathrm{Ism}$ on the direct summands of the factors
\cite[Thm.~3.11(i), p.~154]{IUT3}; the images of possible images under them are possible images. The packet is a
tensor product over the labels of direct sums over the places \cite[Prop.~3.2]{IUT3}, and a morphism of
$\mathcal D^{\vdash}$-prime-strips is a collection of morphisms indexed by the places \cite[Def.~4.1(iv),
p.~96]{IUT1}; so a member of the poly-isomorphism over $\sigma$ maps each summand $k_{\bv}$ to
$k_{\bv\circ\sigma^{-1}}$, sending the factor of label $i$ to that of label $\sigma(i)$ by the map induced by its
component at $i$, and the component at $\bw$ of the image of $X$ is the image of $\pr_{\bv}(X)$. Such members exist,
the members of the capsule being isomorphic \cite[Def.~4.1(iii)]{IUT1}. The log-shells are constructed functorially
from the prime-strips, and the identifications of their $\Q$-spans with the fields $k_v$ form an $\mathrm{Ism}$-orbit
compatible with the log-shells \cite[Prop.~1.2(vi), p.~32]{IUT3}. Hence each member, followed by any element of
(Ind2), maps $\pr_{\bv}(X)$ to $\psi(\sigma_*(\pr_{\bv}X))$ with $\psi_i\in\cA(k_{w_i})$; that is, the images of $X$
are $R(X)$ for maps $R$ of the packet with $\pr_{\bw}\circ R=\psi\circ\sigma_*\circ\pr_{\bv}$.
\end{proof}

\section{The non-containment}\label{s:main}

Let $D_2$ be as in Proposition~\ref{p:datum} and $j=3$, so $\mathbb S^\pm_4=\{0,1,2,3\}$. Let
$\bu=(b_1,b_1,b_2,b_1)$, whose factor of label $3$ lies over $b_1$, let $\theta$ and $x=1\otimes1\otimes1\otimes\theta
\in k_{\bu}$ be as in Lemma~\ref{l:member}, with $\operatorname{ord}(\theta)=9\cdot\tfrac17=\tfrac97$, and let $X\in U_3$
be the element given there, so that $\pr_{\bu}(X)=\chi(x)$. For $\sigma=(2\,3)$ put $\bv=\bu\circ\sigma^{-1}
=(b_1,b_1,b_1,b_2)$ and
\[
y=\sigma_*(x)=1\otimes1\otimes\theta\otimes1\ \in\ k_{b_1}\otimes k_{b_1}\otimes k_{b_1}\otimes k_{b_2}.
\]
By Lemma~\ref{l:transport}, $\pr_{\bv}(U_3)$ contains an element $\psi(\sigma_*(\chi(x)))=\psi'(y)$, where
$\psi'=\otimes_i\psi'_i$ with $\psi'_i=\psi_i\circ\chi_{\sigma^{-1}(i)}\in\cA(k_{v_i})$, since
$\sigma_*\circ\chi=(\otimes_i\chi_{\sigma^{-1}(i)})\circ\sigma_*$.

\begin{theorem}\label{t:main}
For the initial $\Theta$-datum $D_2$, $j=3$ and $\bv=(b_1,b_1,b_1,b_2)$, the projection $\pr_{\bv}(U_3)$ of the
union of possible images of \cite[Cor.~3.12]{IUT3} is not contained in the module $M_{\bv}$ from which
\cite[Thm.~1.10, Step~(v)]{IUT4} obtains its bound for this component. More precisely, no element $\psi(y)$ with
$\psi=\otimes_i\psi_i$, $\psi_i\in\cA(k_{v_i})$, lies in $M_{\bv}$. Consequently the component
$\pr_{\bv}(\overline U_3)$ of the holomorphic hull has log-volume at least $-2\log p$ (possibly $+\infty$), whereas
the bound~\eqref{eq:bound} for it is $-\tfrac{97}{35}\log p$.
\end{theorem}

\begin{proof}
Step~(v) takes $i^\dagger=3$, whose factor lies over $b_2$, so by Table~\ref{tab:D2}
\[
\lambda=9\cdot\tfrac67=\tfrac{54}7,\qquad \dd_I=3\cdot\tfrac{104}{105}+\tfrac{34}{35}=\tfrac{138}{35},\qquad
I^*=\emptyset,
\]
and \eqref{eq:bound} gives $\mulog(M_{\bv})\le(-\tfrac{54}7+\tfrac{138}{35}+1)\log p=-\tfrac{97}{35}\log p$. By (1.3),
$\operatorname{ord}\psi_2(\theta)\le\lceil\tfrac97\rceil=2$ and $\operatorname{ord}\psi_i(1)\le0$ for $i\ne2$; so by
(1.1) every component of $\psi(y)$ has $\operatorname{ord}\le2$, and $\RIt\cdot\psi(y)\supseteq p^2\cdot\RIt$. If
$\psi(y)\in M_{\bv}$, then, $M_{\bv}$ being an $\RIt$-module,
\[
-\tfrac{97}{35}\log p\ \ge\ \mulog(M_{\bv})\ \ge\ \mulog(p^2\cdot\RIt)=-2\log p,
\]
which is false: $-\tfrac{97}{35}<-2$, with margin $\tfrac{27}{35}\log p$. As $\pr_{\bv}(U_3)$ contains an element of
this form, it is not contained in $M_{\bv}$. For the hull: if $U_3$ is not relatively compact, its holomorphic hull
is the whole packet \cite[Rem.~3.9.5(i)]{IUT3}, of log-volume $+\infty$. Otherwise the hull is $\eta\cdot\cO$ for an
element $\eta$ with nonzero components \cite[ibid.]{IUT3}, so $\pr_{\bv}(\overline U_3)$ is an $\RIt$-module
containing $\psi(y)$, hence $p^2\cdot\RIt$, and its log-volume is at least $\mulog(p^2\cdot\RIt)=-2\log p$.
\end{proof}

Thus, for the union of possible images and its holomorphic hull as defined in \cite[Cor.~3.12]{IUT3}, the bound that
Step~(v) obtains for the component at the tuple $\bv$ does not hold. It is noted in \cite[\S6.2, footnote~7]{DH} that
(Ind1) contains permutations between different tensor factors, which the automorphisms used there to estimate the
hull do not include.

The example above concerns the local containment and the local upper bound used in Step~(v) for a single
component. It does not by itself decide the global
inequality obtained later in Theorem~1.10 after summing over all places. No assertion is made here concerning
Corollary~3.12 itself or the abc conjecture.

\subsection*{Use of AI tools}
Generative AI tools (large language models: \mbox{Claude} by \mbox{Anthropic}, and \mbox{ChatGPT} and \mbox{Codex} by \mbox{OpenAI}) were used
extensively in this work, under the author's direction, in
finding the example and its proofs, in the computations and in the writing. The author takes full responsibility for
the contents of this note.

\providecommand{\bysame}{\leavevmode\hbox to3em{\hrulefill}\thinspace}
\providecommand{\MR}{\relax\ifhmode\unskip\space\fi MR }
% \MRhref is called by the amsart/book/proc definition of \MR.
\providecommand{\MRhref}[2]{%
  \href{http://www.ams.org/mathscinet-getitem?mr=#1}{#2}
}
\providecommand{\href}[2]{#2}
\begin{thebibliography}{MFHMP}

\bibitem[CanLift]{CanLift}
Shinichi Mochizuki, \emph{The absolute anabelian geometry of canonical curves},
  Doc. Math. \textbf{Extra Vol.} (2003), 609--640.

\bibitem[DH]{DH}
Taylor Dupuy and Anton Hilado, \emph{Probabilistic {S}zpiro, baby {S}zpiro, and
  explicit {S}zpiro from {M}ochizuki's {C}orollary 3.12}, 2020,
  arXiv:2004.13108v2 [math.NT].

\bibitem[EtTh]{EtTh}
Shinichi Mochizuki, \emph{The \'etale theta function and its
  {F}robenioid-theoretic manifestations}, Publ. Res. Inst. Math. Sci.
  \textbf{45} (2009), no.~1, 227--349.

\bibitem[IUT1]{IUT1}
\bysame, \emph{Inter-universal {T}eichm\"uller theory {I}: construction of
  {H}odge theaters}, Publ. Res. Inst. Math. Sci. \textbf{57} (2021), 3--207,
  page numbers refer to the preprint version of May 2020.

\bibitem[IUT3]{IUT3}
\bysame, \emph{Inter-universal {T}eichm\"uller theory {III}: canonical
  splittings of the log-theta-lattice}, Publ. Res. Inst. Math. Sci. \textbf{57}
  (2021), 403--626, page numbers refer to the preprint version of May 2020.

\bibitem[IUT4]{IUT4}
\bysame, \emph{Inter-universal {T}eichm\"uller theory {IV}: log-volume
  computations and set-theoretic foundations}, Publ. Res. Inst. Math. Sci.
  \textbf{57} (2021), 627--723, page numbers refer to the preprint version of
  April 2020.

\bibitem[MFHMP]{MFHMP}
Shinichi Mochizuki, Ivan Fesenko, Yuichiro Hoshi, Arata Minamide, and Wojciech
  Porowski, \emph{Explicit estimates in inter-universal {T}eichm\"uller
  theory}, Kodai Math. J. \textbf{45} (2022), no.~2, 175--236, numbering as in
  the accepted manuscript of May 2022.

\bibitem[Sil]{Sil}
Joseph~H. Silverman, \emph{The arithmetic of elliptic curves}, second ed.,
  Graduate Texts in Mathematics, vol. 106, Springer, Dordrecht, 2009.

\end{thebibliography}

\end{document}
